% =====================================================================
%  CAPITULO I - REVISION CRITICA DE LA LITERATURA
%  Ficha de tres entradas por estudio: Estudio relevante / Aporte al
%  proyecto / Aporte del proyecto.
%  Todas las cifras citadas provienen de las tablas de los articulos
%  originales. La procedencia exacta de cada numero se indica en una
%  nota al pie, de modo que el jurado pueda verificarla directamente.
% =====================================================================

\chapter{REVISIÓN CRÍTICA DE LA LITERATURA}
\label{cap:revision}

Este capítulo no es un listado de trabajos previos. Es el análisis de los
resultados que publican las cinco fuentes más relevantes de la bibliografía,
hecho con una finalidad concreta: mostrar qué parte del problema de esta tesis
ya está resuelta, qué parte no lo está, y de dónde sale cada uno de los aportes
que el proyecto se compromete a producir. La regla que se sigue es que ningún
aporte se enuncia sin que antes aparezca, en la lectura de un artículo, el
resultado publicado que lo justifica.

Cada estudio se presenta con la misma ficha de tres entradas. La primera,
\textbf{Estudio relevante}, da autores, año y título. La segunda, \textbf{Aporte
al proyecto}, responde tres cosas en este orden: qué hizo el trabajo, qué
resultado obtuvo y por qué ese resultado sirve para esta tesis; es donde se
concentra el análisis crítico de las cifras publicadas. La tercera,
\textbf{Aporte del proyecto}, declara qué se hará distinto, qué se mejorará o en
qué contexto nuevo se aplicará, y a qué objetivo específico queda asociado.

Las cinco publicaciones no se escogieron por cercanía temática. Se escogieron
porque cada una se identifica con un objetivo específico de la tesis y porque
de cada una nace al menos una contribución verificable, según el mapa de la
Tabla~\ref{tab:mapa-estudios}.

\begin{table}[htbp]
\centering
\small
\caption{Correspondencia entre las publicaciones revisadas y los objetivos
específicos. Cada objetivo se sostiene sobre al menos un estudio, y de cada
estudio nace al menos una contribución.}
\label{tab:mapa-estudios}
\begin{tabular}{@{}p{4.6cm}p{4.6cm}p{3.0cm}@{}}
\hline
\textbf{Publicación} & \textbf{Qué aporta a la tesis} & \textbf{Objetivo} \\
\hline
Zhang et al. (2026), HyperNCD \newline (Sección~\ref{sec:hyperncd})
& El método sobre el que se construye y la vara de comparación & OE1 \\
Wu et al. (2025), Sonata \newline (Sección~\ref{sec:sonata})
& La pieza que sustituye al descriptor y la forma de medir representación
& OE1, OE3 \\
Wu et al. (2024), PTv3 \newline (Sección~\ref{sec:ptv3})
& La estructura espacial que crea el problema de integración & OE1 \\
Su et al. (2025), BRDL \newline (Sección~\ref{sec:brdl})
& El procedimiento de calibración bajo desbalance de clases & OE2 \\
Robert et al. (2023), SPT \newline (Sección~\ref{sec:spt})
& La medición aislada del descriptor manual y el reporte de costo & OE3 \\
\hline
\end{tabular}
\end{table}

% =====================================================================
\section[HyperNCD]{Zhang, Wu, Li, Han y Shen (2026): HyperNCD, razonamiento por
hipergrafo}
\label{sec:hyperncd}

\textbf{Estudio relevante.} Zhang, Wu, Li, Han y Shen (2026), \emph{Geometric-Aware
Hypergraph Reasoning for Novel Class Discovery in Point Cloud Segmentation},
publicado en CVPR 2026 \cite{zhang2026hyperncd}. Es el estado del arte del problema y el
método sobre el que esta tesis construye. Se identifica con \textbf{OE1}.

\textbf{Aporte al proyecto.}

\emph{Qué hizo.} Representa las relaciones entre clases con un
\textbf{hipergrafo}. En un grafo común una arista une exactamente dos nodos; en
un hipergrafo una hiperarista puede unir tres, cinco o quince nodos a la vez, lo
que permite expresar afirmaciones del tipo ``estas cinco clases comparten una
misma estructura geométrica'' que un grafo de pares no puede representar. Los
nodos de ese hipergrafo son los prototipos de clase, y esos prototipos se
construyen sobre un descriptor compuesto: por un lado un descriptor semántico
que produce la red supervisada, y por otro un descriptor geométrico
$f_{\text{geo}}$ calculado con una fórmula fija.

\emph{Qué resultado obtuvo.} Sobre SemanticKITTI obtiene entre 11,8 y 47,5 de
mIoU en clases novedosas según la partición, frente a 49,8 del equivalente
supervisado. Sobre SemanticPOSS obtiene entre 22,3 y 47,2, frente a 50,6 del
supervisado.\footnote{Tablas~1 y~2 de \cite{zhang2026hyperncd}. La fila
\emph{Full} de cada tabla es el modelo supervisado y su valor se lee en la
columna \emph{All}.} Más importante que el agregado es su estudio de ablación,
que consiste en apagar una pieza a la vez para ver cuánto aporta cada una. Sobre
la partición~0 de SemanticPOSS, partiendo de una red base MinkowskiUNet-34C
modificada que rinde 23,2 de mIoU promedio, el prototipo consciente de la
geometría aporta $+4{,}2$ puntos; la construcción del hipergrafo aporta $+9{,}6$
sobre la misma base; ambas piezas juntas dan 34,2; y el ajuste dinámico de
hiperaristas eleva el resultado final a 36,5.\footnote{Tabla~3 de
\cite{zhang2026hyperncd}, columna \emph{Overall Avg}.} La
Tabla~\ref{tab:ablacion-hyperncd} recoge esos valores.

\begin{table}[htbp]
\centering
\caption{Estudio de ablación de HyperNCD sobre la partición~0 de SemanticPOSS,
según la Tabla~3 de \cite{zhang2026hyperncd}. La columna de ganancia mide la
diferencia respecto de la red base. Es la evidencia cuantitativa que sostiene la
hipótesis de esta tesis.}
\label{tab:ablacion-hyperncd}
\begin{tabular}{@{}p{6.3cm}cc@{}}
\hline
Configuración & mIoU promedio & Ganancia sobre la base \\
\hline
Red base (MinkowskiUNet-34C modificada) & 23,2 & referencia \\
Solo prototipo consciente de la geometría & 27,4 & $+4{,}2$ \\
Solo hipergrafo & 32,8 & $+9{,}6$ \\
Prototipo geométrico e hipergrafo & 34,2 & $+11{,}0$ \\
Método completo, con hiperaristas dinámicas & 36,5 & $+13{,}3$ \\
\hline
\end{tabular}
\end{table}

\emph{Por qué sirve.} El razonamiento de alto orden entre prototipos funciona, y
la ablación lo cuantifica sin ambigüedad. Pero la misma tabla deja al descubierto
un desequilibrio que los autores no comentan: la pieza geométrica aporta $+4{,}2$
puntos y la pieza de razonamiento aporta $+9{,}6$, es decir, el descriptor
geométrico contribuye menos de la mitad que el mecanismo que lo consume. Conviene
ser preciso con esta lectura, porque de su honestidad depende la validez de la
hipótesis de esta tesis: la ablación \textbf{no} demuestra que el descriptor
geométrico sea inútil; demuestra que aporta poco en relación con el lugar central
que ocupa dentro del método. La hipótesis que se defiende aquí no es que
$f_{\text{geo}}$ no sirva, sino que una representación aprendida puesta en su
lugar debería aportar más.

La limitación de fondo es la naturaleza de ese descriptor. $f_{\text{geo}}$ se
calcula de forma analítica: se toman los $K=15$ vecinos más cercanos de cada
punto, se arma la matriz de covarianza de sus distancias euclidianas, se obtienen
sus autovalores y de ellos se derivan tres cantidades llamadas linealidad,
planaridad y dispersión \cite{zhang2026hyperncd}. Ese cálculo es siempre el
mismo, no depende de los datos y no aprende nada del dominio en el que se aplica.
Además, es exactamente el tipo de señal que la literatura de aprendizaje
auto-supervisado caracteriza como un atajo perjudicial, como se detalla en la
Sección~\ref{sec:sonata}. Una segunda limitación es la dispersión entre
particiones, entre 11,8 y 47,5: los propios autores identifican el desbalance de
clases como el problema que motiva tanto el ajuste dinámico de hiperaristas como
el suavizado de etiquetas con $\epsilon = 0{,}15$. Una tercera, que no está en el
artículo sino en el vacío que deja, es que el trabajo \textbf{no reporta ninguna
cifra de cómputo}: su sección de detalles de implementación declara arquitectura,
optimizador y programa de tasa de aprendizaje, y nada más, sin GPU, sin memoria y
sin tiempo.

De aquí el proyecto toma tres cosas concretas. La arquitectura completa sobre la
que se trabaja, con su código público. Los hiperparámetros publicados, que fijan
el punto de partida de la calibración: $K=15$ vecinos, $M=8$ prototipos vecinos
por hiperarista, $\alpha$ inicializado en $0{,}5$ y $\epsilon = 0{,}15$. Y, sobre
todo, la vara contra la cual se mide el aporte principal: cualquier reemplazo de
$f_{\text{geo}}$ debe compararse contra los $+4{,}2$ puntos que aporta el
descriptor original.

\textbf{Aporte del proyecto.} De este estudio nacen dos contribuciones, ambas
asociadas a \textbf{OE1}.

\begin{enumerate}
\item[\textbf{1.a}] \textbf{Bloque de fusión que sustituye el descriptor
geométrico.} Reemplazar $f_{\text{geo}}$ por el codificador auto-supervisado
congelado de Sonata, $f_{\mathrm{ssl}}$, conservando íntegro el razonamiento por
hipergrafo, de modo que $F_i = [f_{\mathrm{sem}}(x^i);\, f_{\mathrm{ssl}}(x^i)]$
y cualquier variación de la métrica sea atribuible a ese reemplazo y a nada más.
\emph{Resultado verificable:} el híbrido corriendo de extremo a extremo sobre una
sección de control, con las dimensiones de ambas ramas documentadas y sin que el
mIoU de las clases conocidas baje respecto de la línea base.

\item[\textbf{1.c}] \textbf{Auditoría de la implementación de referencia y
corrección de la línea base.} La inspección del código liberado muestra que el
módulo que reproduce la formulación publicada del razonamiento por hipergrafo no
se invoca desde ninguna ruta de entrenamiento ni de evaluación; que las matrices
de peso de sus capas se instancian dentro del cuerpo de la función en cada
llamada, de modo que no se registran entre los parámetros del modelo, no reciben
gradiente y se reinicializan al azar en cada invocación; que el procedimiento
denominado de actualización dinámica reproduce término por término el de
construcción inicial sobre la misma matriz de similitud, por lo que no actualiza
nada; y que el descriptor geométrico efectivamente ejecutado tiene dimensión
seis, porque incorpora el centroide de la región además de las cantidades
derivadas de los autovalores, y no tres como describe el artículo. Estas
observaciones no invalidan el trabajo de referencia, pero obligan a tres
decisiones que sin ellas se habrían tomado mal: la línea base se establece sobre
el código que se ejecuta y no sobre la formulación publicada, la comparación
cuantitativa se hace contra corridas propias bajo condiciones idénticas de datos,
partición y semilla, y la sustitución se aplica sobre el descriptor de dimensión
seis. \emph{Resultado verificable:} la tabla de correspondencia entre artículo e
implementación y las tres decisiones metodológicas que de ella se derivan.
\end{enumerate}

% =====================================================================
\section[Sonata]{Wu et al. (2025): Sonata y el atajo geométrico}
\label{sec:sonata}

\textbf{Estudio relevante.} Wu, DeTone, Frost, Shen, Xie, Yang, Engel, Newcombe,
Zhao y Straub (2025), \emph{Sonata: Self-Supervised Learning of Reliable Point
Representations}, CVPR 2025 \cite{wu2025sonata}. Se identifica con \textbf{OE1}
y, por su metodología de medición, con \textbf{OE3}.

\textbf{Aporte al proyecto.}

\emph{Qué hizo.} Aprendizaje auto-supervisado quiere decir que el modelo aprende
sin etiquetas humanas, resolviendo tareas que se inventa a partir de los propios
datos. El hallazgo central del trabajo es lo que llaman \textbf{atajo
geométrico}: cuando se entrena así sobre nubes de puntos, la representación
tiende a colapsar sobre pistas geométricas de bajo nivel muy fáciles de leer,
como la dirección de la normal de la superficie o la altura del punto, en lugar
de aprender semántica. Los autores argumentan que es un problema propio del 3D y
no de la imagen: en una imagen toda la información está en el contenido de los
píxeles, mientras que en una nube de puntos la geometría entra por las
coordenadas mismas, que los operadores usan de forma directa, y por eso el atajo
es difícil de tapar. Sonata lo ataca oscureciendo la información espacial y
poniendo el énfasis en las características de entrada, dentro de un esquema de
auto-destilación entrenado sobre 140 mil escenas.

\emph{Qué resultado obtuvo.} La evidencia es directa y se mide por \emph{linear
probing}, un procedimiento en el que se congelan todos los pesos del codificador
y se entrena encima únicamente un clasificador lineal, con menos del $0{,}2$\%
de los parámetros. Como el codificador no puede cambiar, lo que se mide es la
calidad de la representación y nada más. Sobre ScanNet el desempeño pasa de 5,6 a
72,5 de mIoU, una mejora que los propios autores describen como $3{,}3$ veces
sobre el estado del arte previo. En S3DIS Área~5 alcanza 72,3 de mIoU con
clasificador lineal, prácticamente a la par de una red supervisada completa, que
obtiene 72,4.\footnote{Tabla~5 de \cite{wu2025sonata}, sobre eficiencia de
parámetros.} La Tabla~\ref{tab:linprobe} recoge la comparación.

\begin{table}[htbp]
\centering
\caption{\emph{Linear probing} sobre ScanNet, conjunto de validación, según la
Tabla~5 de \cite{wu2025sonata}. Todos los métodos congelan el codificador y
entrenan una única capa lineal, con menos del $0{,}2$\% de los parámetros. La
columna mide, por tanto, calidad de la representación y no capacidad del
clasificador.}
\label{tab:linprobe}
\begin{tabular}{lc}
\hline
Método & mIoU (\%) \\
\hline
PointContrast & 5,6 \\
CSC & 12,6 \\
MSC & 21,8 \\
DINOv2, agregado de 2D a 3D & 63,1 \\
\textbf{Sonata} & \textbf{72,5} \\
\hline
\end{tabular}
\end{table}

\emph{Por qué sirve.} Este es el resultado que hace posible la tesis. Lo que
Sonata demuestra es que la señal geométrica de bajo nivel, medida en igualdad de
condiciones, es un mal descriptor, y que existe un codificador público, ya
entrenado y congelable, que la reemplaza sin exigir un preentrenamiento propio.
Sirve además por una segunda razón, de método y no de resultado: el \emph{linear
probing} es la forma de medir calidad de representación de manera aislada,
congelando todo lo demás, y ese principio es el que permite el protocolo de
atribución por etapas que el proyecto necesita para separar lo que aporta la
representación de lo que aporta el razonamiento.

Ahora bien, el trabajo declara tres limitaciones que acotan el riesgo de esta
tesis y que deben enunciarse con la misma precisión:

\begin{enumerate}
\item \textit{Dominio de preentrenamiento.} El modelo principal se preentrena
sobre escenas interiores densas, entre ellas ScanNet, ScanNet++, S3DIS,
ArkitScenes, HM3D, Structured3D y ASE \cite{wu2025sonata}. El levantamiento de
una catedral cae del lado favorable de esa frontera, porque es un interior denso,
pero el cambio de dominio sigue siendo el principal riesgo declarado del
proyecto.

\item \textit{Régimen congelado en exteriores.} En escaneo láser urbano, el
\emph{linear probing} de Sonata queda por debajo de una PTv3 supervisada, con
62,0 frente a 69,1 de mIoU en SemanticKITTI. Hay que registrar, sin embargo, que
con ajuste fino completo Sonata alcanza 72,6 en el mismo conjunto y supera al
supervisado.\footnote{Tabla~8 de \cite{wu2025sonata}.} Es decir, la desventaja es
del régimen congelado y no de la representación en general, y el régimen
congelado es exactamente el que adopta este proyecto, por lo que la limitación
aplica en su forma más restrictiva.

\item \textit{Vocabularios grandes.} Con 200 clases la discriminación cae a 29,3
de mIoU en ScanNet200, frente a 72,5 con 20 clases en ScanNet. Los propios
autores lo señalan como una limitación de la representación aprendida para
distinguir un número grande de categorías. Esto fija de forma directa el límite
superior de clases novedosas simultáneas que el proyecto puede declarar como
rango favorable.
\end{enumerate}

\textbf{Aporte del proyecto.} De este estudio nacen dos contribuciones.

\begin{enumerate}
\item[\textbf{1.a}] \textbf{(continuación) Traslado del hallazgo del atajo
geométrico al régimen de descubrimiento de clases novedosas.} La nocividad del
atajo está demostrada en régimen supervisado y de \emph{linear probing}. Este
trabajo verifica si se mantiene cuando no hay ninguna supervisión sobre las
clases objetivo, que es donde la calidad de la representación pesa más. Es un
resultado incierto, medible y no reportado por nadie. Junto con la sustitución
que describe la Sección~\ref{sec:hyperncd} completa el primer aporte: HyperNCD
aporta el lugar exacto donde se hace el reemplazo y Sonata la pieza que se coloca
en ese lugar y la evidencia de que hacerlo tiene sentido. \emph{Objetivo:} OE1.

\item[\textbf{3.a}] \textbf{Métrica de brecha cerrada $G$ y protocolo de
atribución por etapas.} Las métricas estándar no separan lo que aporta la
representación de lo que aporta el razonamiento: el mIoU final es el mismo número
venga de donde venga la ganancia. Tomando de Sonata el principio del \emph{linear
probing}, este trabajo formula la brecha cerrada $G$ (Ec.~\eqref{eq:G}), que
expresa qué parte de la distancia entre clases conocidas y novedosas se logró
recuperar, y mide también en el punto intermedio del proceso
(Sección~\ref{sec:metricas-hibrido}) para poder atribuir la ganancia a una etapa
concreta. Se incorpora además un análisis de dominancia que identifica qué clase
gobierna cada mejora agregada, de modo que un promedio favorable no oculte una
única clase que arrastra al resto. \emph{Resultado verificable:} $G$ formulada,
la medición intermedia y las proyecciones PCA de los descriptores.
\emph{Objetivo:} OE3.
\end{enumerate}

% =====================================================================
\section[Point Transformer V3]{Wu et al. (2024): Point Transformer V3 y el
problema de integración}
\label{sec:ptv3}

\textbf{Estudio relevante.} Wu, Jiang, Wang, Liu, Liu, Qiao, Ouyang, He y Zhao
(2024), \emph{Point Transformer V3: Simpler, Faster, Stronger}, publicado en
CVPR 2024 \cite{wu2024ptv3}. Es la red sobre la que está construido Sonata: el
codificador que este proyecto emplea congelado es, literalmente, una PTv3 con
pesos auto-supervisados. Se identifica con \textbf{OE1}.

\textbf{Aporte al proyecto.}

\emph{Qué hizo.} Abandonó la vecindad geométrica calculada por vecinos más
cercanos y la sustituyó por una vecindad \emph{serializada}. Los puntos se
cuantizan a una rejilla entera, se ordenan siguiendo una curva de llenado de
espacio, que es una curva que recorre el volumen visitando cada celda una sola
vez y manteniendo juntos los puntos cercanos, la secuencia resultante se corta en
parches contiguos y la atención se calcula únicamente dentro de cada parche. El
trabajo emplea cuatro patrones de recorrido, orden Z, orden de Hilbert y sus dos
variantes transpuestas, y alterna entre ellos para que el campo receptivo no
quede sesgado por la dirección de la curva.

\emph{Qué resultado obtuvo.} La serialización le permite ampliar el campo
receptivo de 16 a 1024 puntos. La justificación de eficiencia es directa: la
consulta de vecinos más cercanos y la codificación posicional relativa ocupaban
en conjunto el 54\,\% del tiempo de paso hacia adelante de la versión anterior, y
PTv3 elimina ambas. Su tabla de eficiencia, medida sobre nuScenes con lote de
tamaño 1 en una sola RTX 4090, da los valores de la
Tabla~\ref{tab:eficiencia-ptv3}.\footnote{Tabla~1 de \cite{wu2024ptv3}.}

\begin{table}[htbp]
\centering
\caption{Eficiencia de PTv3 frente a la red dispersa con la que tendrá que
convivir en el modelo híbrido, según la Tabla~1 de \cite{wu2024ptv3}. Medido
sobre nuScenes, lote de tamaño 1, una sola RTX 4090.}
\label{tab:eficiencia-ptv3}
\begin{tabular}{@{}p{4.2cm}rrrrr@{}}
\hline
 & & \multicolumn{2}{c}{Entrenamiento} & \multicolumn{2}{c}{Inferencia} \\
Método & Parám. & Latencia & Memoria & Latencia & Memoria \\
\hline
MinkUNet, núcleo 3 & 37,9\,M & 163\,ms & 3,3\,G & 48\,ms & 1,7\,G \\
PTv2, parche 16 & 12,8\,M & 213\,ms & 10,3\,G & 146\,ms & 12,3\,G \\
\textbf{PTv3, parche 1024} & 46,2\,M & \textbf{119\,ms} & \textbf{3,3\,G} &
\textbf{44\,ms} & \textbf{1,2\,G} \\
\hline
\end{tabular}
\end{table}

\emph{Por qué sirve.} Sirve porque aquí aparece el problema que este proyecto
tiene que resolver y que ninguno de los dos estudios anteriores plantea, porque
ninguno de los dos intenta juntar las piezas. Las dos ramas del modelo híbrido
parten de las mismas coordenadas, pero definen la vecindad de maneras distintas.
HyperNCD opera sobre una \textbf{rejilla dispersa de Minkowski}: las coordenadas
se cuantizan a celdas y el vecino de un punto son los que caen dentro del núcleo
de la convolución dispersa; la vecindad es adyacencia métrica en el volumen.
Sonata, a través de PTv3, opera sobre una \textbf{secuencia unidimensional}: el
vecino de un punto es el que le quedó al lado dentro de un parche de 1024
elementos de la curva de llenado de espacio. Una curva de este tipo preserva la
localidad de forma aproximada y no exacta, de modo que dos puntos contiguos en el
volumen pueden quedar separados en la secuencia si caen a ambos lados de un giro
de la curva. Cualquier intento de tomar la salida de una rama y alimentarla a la
otra tiene que absorber ese desajuste de forma explícita.

A la incompatibilidad de estructura se suma una segunda, que no es de
arquitectura sino de ingeniería, y que la inspección del código liberado por
ambos grupos deja en evidencia: las dos ramas exigen \textbf{pilas de
dependencias mutuamente excluyentes}, MinkowskiEngine contra CUDA 11.3 y PTv3
contra CUDA 12.4. Esto no es una molestia de instalación. Significa que integrar
las dos implementaciones de referencia dentro de un único proceso no es posible
sin renunciar a una de las dos, y por lo tanto que el punto de fusión no puede
ser un objeto en memoria compartido entre los dos modelos.

La Tabla~\ref{tab:eficiencia-ptv3} aporta además un dato que conviene registrar
porque desactiva una objeción previsible: PTv3 no es la pieza cara del sistema.
Con 46,2\,M de parámetros consume menos memoria de inferencia que MinkUNet, que
tiene 37,9\,M, y es más rápido que ella. El costo del híbrido no proviene de
añadir un transformador, sino de ejecutarlo dentro del bucle de entrenamiento si
no se toma ninguna precaución. Y la propiedad que hace resoluble todo lo anterior
es que el codificador se emplea congelado y solo en inferencia, de modo que su
salida no depende de la época ni del estado del optimizador y puede calcularse
una sola vez por sección.

\textbf{Aporte del proyecto.}

\begin{enumerate}
\item[\textbf{1.b}] \textbf{Puente entre dos estructuras espaciales
incompatibles.} Resolver el emparejamiento entre la secuencia serializada sobre
la que opera PTv3 y la rejilla dispersa de Minkowski sobre la que opera HyperNCD,
mediante \emph{up-cast} de resolución y \emph{cache} de características por
sección. La misma decisión resuelve tres problemas a la vez: el emparejamiento de
vecindades, la incompatibilidad entre pilas de dependencias, que deja de importar
porque las dos ramas nunca coexisten en un mismo proceso, y la repetición del
costo del codificador en cada época, que desaparece. Incluye la formulación
propia de una cota de aceleración que delimita hasta dónde vale la pena
distribuir el cacheo. \emph{Resultado verificable:} el \emph{cache} persistente
por sección, legible desde los dos entornos, y la cota de aceleración.
\emph{Objetivo:} OE1.
\end{enumerate}

% =====================================================================
\section[BRDL]{Su, Chen y Wang (2025): BRDL, calibración bajo desbalance de
clases}
\label{sec:brdl}

\textbf{Estudio relevante.} Su, Chen y Wang (2025), \emph{Balanced Residual
Distillation Learning for 3D Point Cloud Class-Incremental Semantic
Segmentation} \cite{su2025brdl}. Se identifica con \textbf{OE2}.

\textbf{Aporte al proyecto.}

\emph{Qué hizo.} Aborda de forma directa la pregunta que plantea OE2: cómo se
ajusta un sistema que debe incorporar clases sobre las que no tuvo supervisión
completa sin degradar las que ya clasificaba bien. El método combina dos piezas,
una destilación residual que refina el conocimiento previo ampliando la red, y un
aprendizaje balanceado de pseudo-etiquetas que corrige el sesgo hacia las clases
base. El escenario no es el mismo que el de esta tesis y conviene decirlo de
entrada: BRDL trabaja en aprendizaje incremental por clases, donde las categorías
nuevas sí llegan etiquetadas cuando les toca su turno, mientras que en
descubrimiento de clases novedosas nunca lo están.

\emph{Qué resultado obtuvo.} Interesan tres bloques de su evaluación. El primero
es el barrido de un solo hiperparámetro: al mover $\lambda_2$ entre $0{,}1$ y
$1{,}0$ el mIoU recorre el intervalo de 45,1 a 46,28, con el máximo en
$\lambda_2=0{,}2$.\footnote{Tabla~3 de \cite{su2025brdl}, S3DIS, partición S0,
con $|\mathcal{C}_{\text{novel}}|=1$.} El segundo es la grilla conjunta de dos
hiperparámetros acoplados: al barrer $\lambda_2$ y $\lambda_3$ a la vez, el
máximo sube a 46,58 en la combinación $(0{,}2;\,0{,}5)$ y el mínimo se mantiene
en 45,1, de modo que toda la grilla cabe en 1,48 puntos de
mIoU.\footnote{Tabla~4 de \cite{su2025brdl}.} El tercero es la ablación de las
dos estrategias, en la Tabla~\ref{tab:ablacion-brdl}.\footnote{Tabla~6 de
\cite{su2025brdl}, S3DIS, partición S0, con $|\mathcal{C}_{\text{novel}}|=3$,
columna \emph{all}.} A ello se suma un dato de desbalance que el propio artículo
publica sin comentarlo: en el entrenamiento multi-paso una de las clases se queda
en $0{,}00$ de mIoU en los cinco pasos y para los dos métodos comparados, y otra
cae de 34,42 a $0{,}07$.\footnote{Tabla~5 de \cite{su2025brdl}.} Por último,
reporta de forma explícita el costo incremental de su modificación: 142 a 144
segundos por época y 1608 a 1614\,KB de memoria.\footnote{Tabla~8 de
\cite{su2025brdl}.}

\begin{table}[htbp]
\centering
\caption{Ablación de las dos estrategias de BRDL sobre S3DIS, partición S0,
según la Tabla~6 de \cite{su2025brdl}. La fila intermedia es la que importa para
esta tesis: una de las dos estrategias, aplicada por separado, deja el resultado
por debajo de no aplicar ninguna.}
\label{tab:ablacion-brdl}
\begin{tabular}{@{}p{6.3cm}cc@{}}
\hline
Configuración & mIoU (\emph{all}) & Diferencia \\
\hline
Sin ninguna de las dos estrategias & 45,19 & referencia \\
Solo una de las dos estrategias & 44,90 & $-0{,}29$ \\
Solo la otra estrategia & 46,83 & $+1{,}64$ \\
Las dos estrategias juntas & 47,26 & $+2{,}07$ \\
\hline
\end{tabular}
\end{table}

\emph{Por qué sirve.} Por tres razones y con una advertencia. La primera es
metodológica: es el único de los trabajos revisados que publica una grilla
conjunta de dos parámetros acoplados y que la lee como tal, concluyendo que
$\lambda_2$ domina el resultado y que $\lambda_3$ solo introduce ajustes finos
que amplifican o amortiguan el efecto según el valor que tome $\lambda_2$. El
efecto no es separable, y una calibración parámetro por parámetro llega a un
óptimo distinto del que encuentra la grilla conjunta. La
Tabla~\ref{tab:ablacion-brdl} refuerza el mismo punto desde el lado de los
componentes: activar una sola de las dos piezas empeora el resultado respecto de
no activar ninguna, y solo la combinación gana.

La segunda es de magnitud, y sirve para no prometer de más. Toda la grilla de
hiperparámetros mueve 1,48 puntos de mIoU, mientras que activar o desactivar los
componentes mueve 2,07. La calibración es necesaria, pero no es donde está la
ganancia grande, y esa cifra fija la expectativa realista de OE2.

La tercera es sobre el desbalance, y es la más dura: el fracaso por falta de
soporte no es gradual sino total, con una clase en cero en los cinco pasos y para
los dos métodos. Ningún valor de hiperparámetro rescata una clase cuyo soporte es
insuficiente. Eso coincide con el régimen que esta tesis declara desfavorable, el
de clases novedosas por debajo del 0,5\% de los puntos de la sección, y explica
por qué HyperNCD necesita a la vez el ajuste dinámico de hiperaristas y el
suavizado de etiquetas $\epsilon$ (Sección~\ref{sec:hyperncd}).

La advertencia es la diferencia de escenario ya señalada. Las cifras de BRDL no
se pueden trasladar; lo que se traslada es el procedimiento y el diagnóstico.

\textbf{Aporte del proyecto.} De este estudio nacen las tres contribuciones de
\textbf{OE2}.

\begin{enumerate}
\item[\textbf{2.a}] \textbf{Normalización de escala entre dos ramas
heterogéneas.} La afinidad entre prototipos combina una componente geométrica y
una semántica según $S_b=\alpha\,S_g+(1-\alpha)\,S_s$. Mientras $S_g$ se calcula
sobre un descriptor de seis componentes acotadas derivadas de autovalores, el
valor publicado $\alpha=0{,}5$ significa efectivamente mitad y mitad. En cuanto
$S_g$ pasa a calcularse sobre un vector de cientos de dimensiones producido por
un codificador congelado, cuyo rango no está normalizado contra el de la rama
semántica, ese mismo valor deja de significar eso. Decidir cómo se hacen
comparables las dos similitudes es una precondición de la calibración y no un
detalle de implementación, y el régimen congelado agrava el problema porque
impide que la rama aprendida se reescale sola durante el entrenamiento. En BRDL
los dos parámetros pertenecen al mismo método y comparten escala, de modo que
este caso queda fuera de lo que ese trabajo cubre. \emph{Resultado verificable:}
el procedimiento de normalización adoptado y su efecto medido sobre la
distribución de $S_b$ antes y después de aplicarlo.

\item[\textbf{2.b}] \textbf{Calibración conjunta de $(\alpha, M, \epsilon)$ y
verificación de no-aditividad.} Determinar los tres valores con una grilla
conjunta y curvas de sensibilidad sobre una sección de referencia, en lugar de un
barrido por parámetro, y verificar si la no-aditividad que BRDL documenta entre
sus dos estrategias aparece también entre la representación y el razonamiento del
híbrido. \emph{Resultado verificable:} la grilla conjunta, las tres curvas de
sensibilidad y la configuración fijada.

\item[\textbf{2.c}] \textbf{Verificación de que la configuración se sostiene
fuera de donde se calibró.} La configuración se fija sobre una única sección de
referencia y se aplica \textbf{sin reajuste} a las cuatro restantes, reportando
la degradación observada. Es lo que separa haber calibrado de haber ajustado
sobre aquello que se iba a medir, y el estudio de referencia muestra por qué
importa: al cambiar únicamente el orden de presentación de las clases, los
métodos incrementales existentes se degradan de forma significativa, mientras que
el entrenamiento conjunto apenas se mueve. \emph{Resultado verificable:} la
métrica por sección bajo configuración fija y su diferencia respecto de la
sección de calibración.
\end{enumerate}

% =====================================================================
\section[Superpoint Transformer]{Robert, Raguet y Landrieu (2023): Superpoint
Transformer y el valor medido del descriptor manual}
\label{sec:spt}

\textbf{Estudio relevante.} Robert, Raguet y Landrieu (2023), \emph{Efficient 3D
Semantic Segmentation with Superpoint Transformer}, ICCV 2023
\cite{robert2023spt}. Se identifica con \textbf{OE3}.

\textbf{Aporte al proyecto.}

\emph{Qué hizo.} Es una arquitectura de segmentación semántica que agrupa la nube
en \emph{superpuntos}, es decir en regiones que se comportan como una unidad, y
razona sobre las relaciones entre esas regiones en lugar de sobre puntos
individuales. Se incluye en esta revisión por una razón puntual y no por su
arquitectura: es el único trabajo revisado que mide de forma aislada, y sobre
tres conjuntos de datos distintos, exactamente la operación que esta tesis
ejecuta, que es quitar el descriptor geométrico calculado a mano.

\emph{Qué resultado obtuvo.} Alcanza 76,0 de mIoU en S3DIS con validación de seis
pliegues, 63,5 en KITTI-360 y 79,6 en DALES, con 212 mil parámetros, hasta 200
veces menos que los métodos comparables, y se entrena en una sola GPU en 1,9
horas frente a las 216,4 horas de GPU de MinkowskiNet.\footnote{Tablas~2 y~3 de
\cite{robert2023spt}.} El resultado que importa aquí es otro: en su estudio de
ablación, la fila \emph{sin descriptores manuales} cuesta $-0{,}7$ de mIoU en
S3DIS, $-4{,}1$ en KITTI-360 y $-1{,}4$ en DALES.\footnote{Tabla~4 de
\cite{robert2023spt}, fila a).}

\emph{Por qué sirve.} Ese único renglón obliga a enunciar con más precisión la
hipótesis del capítulo. Puesto junto a la ablación de HyperNCD
(Tabla~\ref{tab:ablacion-hyperncd}) produce la comparación de la
Tabla~\ref{tab:canal-geo}, que hasta donde alcanza esta revisión no ha sido
reunida en ningún trabajo previo.

\begin{table}[htbp]
\centering
\caption{Cuánto vale el canal geométrico calculado a mano, según dos grupos
independientes y tres conjuntos de datos. Las cifras de \cite{robert2023spt} son
pérdidas al quitar el descriptor; la de \cite{zhang2026hyperncd} es la ganancia
al agregarlo, de modo que las operaciones son inversas y no idénticas, pero la
magnitud es comparable. La celda final es la que este proyecto llena.}
\label{tab:canal-geo}
\begin{tabular}{@{}p{3.2cm}p{3.8cm}p{2.2cm}c@{}}
\hline
Conjunto de datos & Régimen & Trabajo & Efecto (mIoU) \\
\hline
S3DIS, 6 pliegues & Interior denso & \cite{robert2023spt} & $-0{,}7$ \\
DALES & Aéreo & \cite{robert2023spt} & $-1{,}4$ \\
KITTI-360 & Exterior disperso & \cite{robert2023spt} & $-4{,}1$ \\
SemanticPOSS, part.~0 & Exterior disperso & \cite{zhang2026hyperncd} & $+4{,}2$ \\
Catedral de Lima & Interior denso, NCD & Este trabajo & por medir \\
\hline
\end{tabular}
\end{table}

Dos observaciones se desprenden de ahí. La primera es la convergencia: dos grupos
independientes, con dos arquitecturas que no tienen nada en común, miden el canal
geométrico manual en alrededor de cuatro puntos de mIoU cuando el escaneo es
exterior y disperso. La segunda, que es la que afecta directamente a este
proyecto, es que en interior denso ese mismo canal vale $-0{,}7$, casi seis veces
menos. El objeto de estudio de esta tesis, el levantamiento interior de una
catedral, cae del lado de S3DIS y no del lado de KITTI-360. Eso corta en las dos
direcciones y así hay que declararlo: la vara de $+4{,}2$ puntos que fija la
Sección~\ref{sec:hyperncd} es una vara medida en exterior disperso, de modo que
en el dominio de este trabajo hay menos que superar, pero también menos margen
disponible para que la sustitución muestre una ganancia distinguible del ruido
entre corridas.

Hay que registrar además que los autores de SPT leen su propio resultado en
dirección contraria a la hipótesis de esta tesis: escriben que su observación
coincide con la de otros trabajos que señalan el impacto positivo de descriptores
manuales bien diseñados en modelos pequeños. No conviene disimular esa lectura.
SPT no afirma que el descriptor manual sobre; afirma que ayuda, y que en modelos
de pocos parámetros ayuda de forma consistente. La hipótesis defendida aquí no la
contradice, porque el proyecto no propone eliminar el canal sino sustituirlo por
otro más informativo. Lo que SPT fija es la cota que el reemplazo debe superar, y
la fija por dominio en lugar de en agregado.

La tercera contribución de este trabajo a la revisión es de método de reporte.
SPT informa el costo en columnas separadas, preproceso en minutos, entrenamiento
en horas de GPU, inferencia en segundos y número de parámetros, y demuestra que
esa desagregación cambia la conclusión: un método puede ganar en mIoU y perder
por completo cuando se mira el preproceso, como ocurre con los 89,9 frente a 12,4
minutos que compara su Tabla~3. Un sistema híbrido que ejecuta un codificador
congelado adicional es exactamente el caso en el que ese reporte desagregado no
puede omitirse, y más todavía cuando el estado del arte del problema no publica
ninguna cifra de cómputo (Sección~\ref{sec:hyperncd}).

\textbf{Aporte del proyecto.} De este estudio nacen dos de las tres
contribuciones de \textbf{OE3}.

\begin{enumerate}
\item[\textbf{3.b}] \textbf{Medición del canal geométrico en la combinación de
dominio y régimen que nadie ha reportado.} Llenar la celda final de la
Tabla~\ref{tab:canal-geo}: interior denso patrimonial bajo descubrimiento de
clases novedosas, que es exactamente el objeto de estudio de esta tesis. Medirla
determina de qué lado cae el balance entre tener menos vara que superar y tener
menos margen disponible. \emph{Resultado verificable:} la
Tabla~\ref{tab:canal-geo} completa, con su última fila medida, y el estudio de
ablación de las tres variantes del descriptor, que son la geométrica manual, la
auto-supervisada y la concatenación de ambas.

\item[\textbf{3.c}] \textbf{Reporte desagregado del costo, que la literatura de
descubrimiento de clases novedosas no publica.} Reportar preproceso, horas de
GPU, memoria y parámetros entrenados por separado, sobre hardware declarado y
verificado, y comparar además contra el enfoque supervisado de vocabulario
cerrado aplicado previamente al mismo tipo de edificio \cite{tesisutec2024}.
\emph{Resultado verificable:} la tabla de costo desagregado, con las mediciones
directas que el procedimiento de cacheo registra por sección.
\end{enumerate}

% =====================================================================
\section[Conclusión: aportes del proyecto]{Conclusión de la revisión crítica}
\label{sec:aporte}

Los estudios previos resuelven el problema \textbf{solo de forma parcial}. Existe
una formulación establecida del descubrimiento de clases novedosas y un método de
razonamiento entre clases que funciona, pero su descriptor geométrico aporta menos
de la mitad que el mecanismo que lo consume (Sección~\ref{sec:hyperncd}). Existe
una representación de puntos de calidad demostrada, pero medida en régimen
supervisado y de vocabulario cerrado (Sección~\ref{sec:sonata}). Existe una
caracterización precisa de la estructura espacial sobre la que esa representación
opera, y precisamente por eso se sabe que no coincide con la del método anterior
(Sección~\ref{sec:ptv3}). Existe un procedimiento publicado para calibrar
mecanismos acoplados, pero sobre parámetros de un mismo método y de una misma
escala (Sección~\ref{sec:brdl}). Y existe una medición aislada de cuánto vale el
descriptor manual, pero en tres dominios que no incluyen el de esta tesis
(Sección~\ref{sec:spt}). Ninguno de esos trabajos se ha evaluado sobre patrimonio
construido, y nadie ha juntado las piezas.

Sobre esa base, el proyecto declara \textbf{tres aportes principales, uno por
objetivo específico}, descompuestos en nueve contribuciones que ya se enunciaron
en la sección de la que nace cada una. El inventario completo, con su origen y su
resultado verificable, está en la Tabla~\ref{tab:obj-vs-revision}.

\medskip
\noindent\textbf{Aporte 1 --- Integración de Sonata en el módulo de prototipos
(OE1).} Sustituir el descriptor geométrico analítico por el codificador
auto-supervisado congelado, conservando íntegro el razonamiento por hipergrafo,
resolver el puente entre las dos estructuras espaciales que hace viable esa
sustitución, y corregir la base de comparación a partir de la auditoría del
código liberado. Nace de las Secciones~\ref{sec:hyperncd}, \ref{sec:sonata}
y~\ref{sec:ptv3}.

\medskip
\noindent\textbf{Aporte 2 --- Calibración del híbrido (OE2).} Hacer comparables
dos ramas de escalas distintas, calibrar de forma conjunta los tres parámetros
que gobiernan la afinidad entre prototipos, y verificar que la configuración
elegida se sostiene fuera de la sección donde se calibró. Nace de la
Sección~\ref{sec:brdl}.

\medskip
\noindent\textbf{Aporte 3 --- Cuantificación de la mejora y delimitación del
rango de operación (OE3).} Formular una métrica que separe la contribución de la
representación de la del razonamiento, medir el valor del canal geométrico en la
única combinación de dominio y régimen que ningún trabajo ha reportado, y
reportar el costo de forma desagregada sobre hardware verificado. Nace de las
Secciones~\ref{sec:sonata} y~\ref{sec:spt}.

\begin{table}[htbp]
\centering
\small
\caption{Inventario de aportes. Tres aportes principales, uno por objetivo
específico, descompuestos en nueve contribuciones. Cada una nace de un estudio
revisado y termina en un resultado propio que se puede verificar.}
\label{tab:obj-vs-revision}
\begin{tabular}{@{}p{0.7cm}p{4.3cm}p{2.8cm}p{4.7cm}@{}}
\hline
 & \textbf{Contribución} & \textbf{Nace de} & \textbf{Resultado verificable} \\
\hline
\multicolumn{4}{@{}l}{\textbf{Aporte 1 --- Integración de Sonata en los prototipos (OE1)}}\\
1.a & Bloque de fusión que sustituye $f_{\text{geo}}$ por $f_{\mathrm{ssl}}$
    & HyperNCD (\S\ref{sec:hyperncd}), Sonata (\S\ref{sec:sonata})
    & Híbrido de extremo a extremo sin caída en clases conocidas \\
1.b & Puente entre serialización y rejilla dispersa
    & PTv3 (\S\ref{sec:ptv3})
    & \emph{Cache} por sección y cota de aceleración propia \\
1.c & Auditoría del código y corrección de la línea base
    & HyperNCD (\S\ref{sec:hyperncd}) y su repositorio
    & Tabla artículo\,--\,implementación y tres decisiones metodológicas \\
\hline
\multicolumn{4}{@{}l}{\textbf{Aporte 2 --- Calibración del híbrido (OE2)}}\\
2.a & Normalización de escala entre ramas heterogéneas
    & BRDL (\S\ref{sec:brdl}), Sonata (\S\ref{sec:sonata})
    & Distribución de $S_b$ antes y después \\
2.b & Calibración conjunta de $(\alpha,M,\epsilon)$ y no-aditividad
    & BRDL (\S\ref{sec:brdl})
    & Grilla conjunta, curvas de sensibilidad, configuración fijada \\
2.c & Robustez de la configuración entre secciones
    & BRDL (\S\ref{sec:brdl})
    & Métrica por sección sin reajuste \\
\hline
\multicolumn{4}{@{}l}{\textbf{Aporte 3 --- Cuantificación y rango de operación (OE3)}}\\
3.a & Brecha cerrada $G$ y atribución por etapas
    & Sonata (\S\ref{sec:sonata})
    & $G$ formulada, medición intermedia, PCA \\
3.b & Valor del canal geométrico en interior denso bajo NCD
    & SPT (\S\ref{sec:spt}), HyperNCD (\S\ref{sec:hyperncd})
    & Tabla~\ref{tab:canal-geo} completa y ablación de tres variantes \\
3.c & Reporte desagregado del costo
    & SPT (\S\ref{sec:spt}) y el vacío de HyperNCD
    & Tabla de costo por etapa, con mediciones directas \\
\hline
\end{tabular}
\end{table}

La contribución al estado del arte no es, por tanto, una arquitectura nueva. Es
establecer si un hallazgo demostrado en un régimen se transfiere a otro donde
resulta más crítico, hacerlo sobre un objeto de estudio en el que el costo del
trabajo manual es real y medible, y dejar documentado con qué base de comparación
y a qué costo se hizo. Cada punto de mIoU que se gana en las clases novedosas es
trabajo de corrección que el especialista deja de hacer, y por eso el valor del
proyecto no está en igualar la precisión del trabajo manual sino en automatizarlo
hasta un nivel que resulte útil.
